\begin{statement}[C.8.9; Refuted sign and normalization, with the corrected current identity proved]
On a local period chart write $\mu=a\tau+b$, with $a,b$ real,
and put $T=\Im\tau$. The positive Betti measure is
\[
 m_P=db\wedge da
 =\frac{|\mu_w-a\tau_w|^2}{T}\,dA_w.
\]
Thus the form $da\wedge db$ in Lead 15(b) has the wrong sign.
The globally single-valued function $L(t)=\log N(P_t)$ satisfies
on $\mathbb P^1(\C)$
\begin{equation}\label{eq:ca-current}
 dd^cL=\frac1{12}\delta_\infty-m_P,\qquad
               \int_{\mathbb P^1}m_P=\frac1{12}.
\end{equation}
There are no atoms of $m_P$ at $0$, $1$ or $\infty$.
Let $T_0=\sqrt3/2$, $\rho=(1+i\sqrt3)/2$, and
$p=\wp'_\rho((2-\rho)/3)$. The density in the ordinary
$t$ coordinate satisfies
\begin{align}
 \frac{dm_P}{dA_t}
 &=\frac{|t|^{-2/3}}{64T_0|p|^{2/3}}
                      \bigl(1+O(|t|^{1/3})\bigr),
                         &&t\longrightarrow0,\label{eq:ca-density0}\\
 \frac{dm_P}{dA_t}
 &=\frac{|t-1|^{-3/2}}{24\pi^2\sqrt{E_4(i)}}
                      \bigl(1+O(|t-1|^{1/2})\bigr),
                         &&t\longrightarrow1.\label{eq:ca-density1}
\end{align}
\end{statement}
\begin{proof}
On a chart avoiding a zero of the section, the theta quotient
is holomorphic and nonzero. Its log modulus is pluriharmonic.
Only $-\pi(\Im\mu)^2/T$ contributes to the curvature.
Direct differentiation of this real quotient gives
\[
 \partial_w\partial_{\bar w}\frac{(\Im\mu)^2}{T}
       =\frac{|\mu_w-a\tau_w|^2}{2T}.
\]
Also
\[
 a_w=\frac{\mu_w-a\tau_w}{2iT},\qquad
 b_w=-\bar\tau\frac{\mu_w-a\tau_w}{2iT}.
\]
Taking the wedge product proves the formula for $m_P$ and
$dd^cL=-m_P$ on the chart. The transformation of $(a,b)$
under an integral change of period basis has determinant one,
so the measure patches. Translation of $\mu$ by a locally
constant integral period changes neither the formula nor
the measure. The formula remains valid on a root chart
at an elliptic modulus.

At $t=0$ the expansion proved in C.3 is
\[
 L(t)=\frac14\log3+\frac7{192}\Re t
       -B|t|^{4/3}+O(|t|^{5/3}),\qquad
 B=\frac{9\pi}{512T_0|p|^{2/3}}.
\]
The expansion is a convergent real-analytic expansion on a
root chart, so it can be differentiated termwise with the
indicated orders. The Euclidean identity
$\Delta_{\R^2}|t|^\alpha=\alpha^2|t|^{\alpha-2}$ gives
\eqref{eq:ca-density0}.

At $t=1$ put $h=\mu-(1-\tau)/2$. The checked K3 expansion
and the modular $j$ expansion give
\[
 t-1=A_1h^4(1+O(h^2)),\qquad
 A_1=\frac94\pi^4E_4(i)>0,\qquad
 L=\log\sqrt2-\frac\pi2|h|^2+O(|h|^4).
\]
Therefore
$L=\log\sqrt2-\pi|t-1|^{1/2}/(2\sqrt{A_1})+O(|t-1|)$.
Termwise differentiation on the root chart proves
\eqref{eq:ca-density1}. Both densities are locally integrable,
and the continuous potentials at $0$ and $1$ have no atoms.

For infinity take $u=1/t$, and set $\ell=\log(5+2\sqrt6)$.
The cusp branch is
$\mu=1/2-i\ell/(2\pi)+O(q)$, $q=e^{2\pi i\tau}$.
The theta and eta products, with the nonzero factor
$|2\sin\pi\mu|=2\sqrt3+O(q)$, give
\begin{equation}\label{eq:ca-cuspexp}
 L=\frac1{12}\log|u|+\log(2\sqrt3)-\frac1{12}\log1728
          -\frac{\ell^2}{2\log(1728/|u|)}
          +O\left(\frac{|u|}{|\log|u||}+|u|\right).
\end{equation}
Here $q=u/1728+O(u^2)$; the estimate follows directly by
bounding the convergent products. The nonlogarithmic
correction is continuous at $u=0$ and has no point mass.
Equivalently its Betti density is
\[
 \frac{dm_P}{dA_u}
  =\frac{\ell^2}{2\pi|u|^2\log(1728/|u|)^3}
       \bigl(1+o(1)\bigr),
\]
which is integrable. The single logarithmic term in
\eqref{eq:ca-cuspexp} contributes $\delta_\infty/12$.
This proves the global current identity; pairing it with
the constant function one gives the total mass.
\end{proof}
